%% kapa: exam=1 round=32 year=2021 period=1 form=A subject=economics subject_ko=경제학원론 seq=30 number=70
%% source: Q-Net 감정평가사 기출문제 (cst00309) artlSeq=5208397 file=extracted/1008/5208397/A형/2021년도 제32회 감정평가사 1차 1교시 A형.pdf
%% conversion: kordoc 4.24.51 -> tools/gen1.py (auto); review=proofread-sonnet; flags=none
\begin{question}{70}
표는 기업 甲과 乙로만 구성된 A국의 연간 국내 생산과 분배를 나타낸다. 이에 관한 설명으로 옳지 않은 것은?
\par{\small\centering
\begin{tabular}{|l|c|c|}\hline
\multicolumn{1}{|c|}{항목} & 甲 & 乙 \\ \hline
매출액 & 400 & 900 \\ \hline
중간투입액 & 0 & 400 \\ \hline
임금 & 250 & 300 \\ \hline
이자 & 0 & 50 \\ \hline
임대료 & 100 & 100 \\ \hline
이윤 & (\hspace{1em}) & (\hspace{1em}) \\ \hline
요소소득에 대한 총지출 & (\hspace{1em}) & (\hspace{1em}) \\ \hline
부가가치 & (\hspace{1em}) & (\hspace{1em}) \\ \hline
\end{tabular}\par}
\begin{choices}
  \item 기업 甲의 요소소득에 대한 총지출은 400이다.
  \item 기업 甲의 부가가치는 400이다.
  \item 기업 甲의 이윤은 기업 乙의 이윤과 같다.
  \item A국의 임금, 이자, 임대료, 이윤에 대한 총지출은 900이다.
  \item A국의 국내총생산은 기업 甲와 기업 乙의 매출액 합계에서 요소소득에 대한 총지출을 뺀 것과 같다.
\end{choices}
\end{question}
\answer{5}
